% ===========================================================================
%  Chapitre 13 — Probabilité
%  PE 2026, composante TRAITEMENT DES DONNÉES ET PROBABILITÉS
% ===========================================================================
\bmtchapitre{Probabilité}
  {Employer le vocabulaire probabiliste, reconnaître une situation
   d'équiprobabilité, calculer la probabilité d'un événement dans les trois
   modèles de tirage, et utiliser l'événement contraire.}
  {Traitement des données \textbullet\ environ 8 heures}

Le chapitre précédent a appris à compter ; celui-ci apprend à conclure. La
formule tient en une ligne — nombre de cas favorables divisé par nombre de
cas possibles — et tout le travail consiste à remplir correctement les deux
étages de la fraction.

C'est l'un des deux exercices de cinq points de l'épreuve, et le plus
régulier de tous : une urne, un sac ou une trousse ; des boules, des jetons,
des billets ou des stylos ; deux ou trois tirages ; quatre ou cinq
événements à évaluer. Vingt-sept sessions se ressemblent ainsi.

\begin{remarque}[ce qui n'est pas au programme]
Les \textbf{probabilités conditionnelles}, l'\textbf{indépendance} de deux
événements et les \textbf{variables aléatoires} sont explicitement hors
programme en série littéraire. Aucun sujet ne les a jamais demandés, et ce
livre ne les traite pas. L'univers est toujours \emph{fini} et les issues
toujours \emph{équiprobables}.
\end{remarque}

% ---------------------------------------------------------------------------
\begin{revision}
Le chapitre 12 suffit.

\begin{enumerate}[label=\textbf{\arabic*.}, leftmargin=*, itemsep=0.35em]
  \item Calculer $\combi{2}{6}$ et $\arrang{2}{6}$.
  \item Une urne contient $4$ boules rouges et $6$ vertes : combien de
        tirages simultanés de deux boules ?
  \item Simplifier la fraction $\dfrac{36}{120}$.
  \item Que vaut $1-\dfrac{7}{15}$ ?
  \item Un tirage « successivement et avec remise » de $3$ objets parmi
        $8$ : combien de cas possibles ?
  \item Donner un nombre compris entre $0$ et $1$ sous forme de fraction
        irréductible.
\end{enumerate}
\end{revision}

\begin{automatismes}
Sans calculatrice.

\begin{enumerate}[label=\textbf{\alph*)}, leftmargin=*, itemsep=0.25em]
  \item Probabilité d'obtenir « pile » avec une pièce équilibrée
  \item Probabilité d'obtenir un $6$ avec un dé à six faces
  \item Probabilité d'obtenir un numéro pair avec ce dé
  \item $1-\dfrac{1}{4}$
  \item $\dfrac{10}{66}$ sous forme irréductible
  \item Probabilité de l'événement certain
  \item Probabilité de l'événement impossible
  \item Si $\proba{A} = 0{,}3$, que vaut $\proba{\evtc{A}}$ ?
\end{enumerate}
\end{automatismes}

\begin{corrige}[automatismes]
a) $\frac12$ \quad b) $\frac16$ \quad c) $\frac12$ \quad d) $\frac34$ \quad
e) $\frac{5}{33}$ \quad f) $1$ \quad g) $0$ \quad h) $0{,}7$.
\end{corrige}

% ===========================================================================
\section{Le vocabulaire}

\begin{definition}[expérience aléatoire, univers]
Une \textbf{expérience aléatoire} est une expérience dont on connaît tous
les résultats possibles sans pouvoir prévoir celui qui se produira.
L'ensemble de ces résultats s'appelle l'\textbf{univers}, noté $\Omega$ ;
chacun d'eux est une \textbf{éventualité}.
\end{definition}

\begin{definition}[événement]
Un \textbf{événement} est une partie de l'univers. On dit qu'il est
\textbf{réalisé} si le résultat de l'expérience lui appartient.
\begin{itemize}[leftmargin=1.3em, itemsep=0.12em]
  \item L'événement $\Omega$ est \textbf{certain}, l'ensemble vide est
        \textbf{impossible}.
  \item Un événement réduit à une seule éventualité est dit
        \textbf{élémentaire}.
  \item L'événement \textbf{contraire} de $A$, noté $\evtc{A}$, est réalisé
        exactement quand $A$ ne l'est pas.
  \item $A \cup B$ est réalisé si $A$ \emph{ou} $B$ l'est ; $A \cap B$ si
        $A$ \emph{et} $B$ le sont. Si $A \cap B = \varnothing$, les deux
        événements sont dits \textbf{incompatibles}.
\end{itemize}
\end{definition}

\begin{visualisation}[deux événements dans un univers]
\begin{center}
\begin{tikzpicture}[scale=0.95]
  \draw[bmtGris, line width=0.8pt, rounded corners=4pt]
    (0,0) rectangle (7.4,3.5);
  \node[font=\footnotesize, text=bmtGris, anchor=north west]
    at (0.12,3.42) {$\Omega$};
  \begin{scope}
    \clip (2.5,1.75) circle (1.3);
    \fill[bmtRaphia!30!bmtPapier] (4.4,1.75) circle (1.3);
  \end{scope}
  \draw[bmtIndigo, line width=1.1pt] (2.5,1.75) circle (1.3);
  \draw[bmtLaterite, line width=1.1pt] (4.4,1.75) circle (1.3);
  \node[bmtIndigo, font=\footnotesize] at (1.75,1.75) {$A$};
  \node[bmtLaterite, font=\footnotesize] at (5.15,1.75) {$B$};
  \node[font=\scriptsize, text=bmtRaphia] at (3.45,1.05) {$A \cap B$};
  \node[font=\scriptsize, text=bmtGris, anchor=west] at (6.1,0.42)
    {$\evtc{A}$};
\end{tikzpicture}
\end{center}

Tout ce qui est à l'extérieur du disque $A$ constitue son contraire
$\evtc{A}$. La zone commune aux deux disques est $A \cap B$ ; lorsque les
disques ne se rencontrent pas, les événements sont incompatibles, et c'est
seulement dans ce cas que les probabilités s'additionnent simplement.
\end{visualisation}

% ===========================================================================
\section{Probabilité d'un événement}

\begin{definition}[cas d'équiprobabilité]
Lorsque toutes les éventualités ont la même chance de se produire — on dit
qu'il y a \textbf{équiprobabilité} —, la probabilité d'un événement $A$ est
\[
  \proba{A} = \frac{\text{nombre de cas favorables}}
                   {\text{nombre de cas possibles}}
            = \frac{\card{A}}{\card{\Omega}} .
\]
\end{definition}

\begin{remarque}
Les sujets signalent l'équiprobabilité par une formule rituelle :
« indiscernables au toucher », « chaque boule a la même probabilité d'être
tirée », « on suppose les tirages équiprobables ». Dès que vous la lisez,
vous savez que la formule ci-dessus s'applique.
\end{remarque}

\begin{propriete}[règles de calcul]
\[
  0 \leq \proba{A} \leq 1, \qquad \proba{\Omega} = 1, \qquad
  \proba{\varnothing} = 0,
\]
\[
  \proba{\evtc{A}} = 1-\proba{A}, \qquad
  \proba{A \cup B} = \proba{A}+\proba{B}-\proba{A \cap B} .
\]
Si $A$ et $B$ sont incompatibles, cette dernière égalité se réduit à
$\proba{A \cup B} = \proba{A}+\proba{B}$.
\end{propriete}

\begin{avertissement}
Une probabilité est un nombre compris entre $0$ et $1$. Un résultat comme
$\frac{7}{5}$ ou $-0{,}2$ signale une erreur de décompte, et il faut
reprendre le calcul plutôt que d'écrire la réponse. Les sujets demandent en
outre, presque toujours, les résultats \textbf{sous forme de fraction
irréductible}.
\end{avertissement}

\begin{methode}{« au moins », « au plus » : passer par le contraire}
\begin{itemize}[leftmargin=1.3em, itemsep=0.2em]
  \item « \textbf{Au moins un} » a pour contraire « \textbf{aucun} », qui se
        compte en une seule ligne. C'est presque toujours le chemin le plus
        court.
  \item « \textbf{Au plus un} » signifie « aucun \emph{ou} exactement un » :
        on additionne les deux décomptes, les cas étant incompatibles.
  \item « \textbf{Exactement} $k$ » se compte directement, en choisissant
        les $k$ objets d'une catégorie et les autres dans le reste.
\end{itemize}
\end{methode}

\begin{exemple}[une urne, quatre événements]
Une urne contient $5$ jetons indiscernables : $3$ noirs et $2$ blancs. On
tire simultanément $3$ jetons. Il y a $\combi{3}{5} = 10$ cas possibles.

\textbf{$A$ : « trois jetons de même couleur ».} Seuls les noirs sont assez
nombreux : $\combi{3}{3} = 1$ cas, donc $\proba{A} = \frac{1}{10}$.

\textbf{$B$ : « exactement deux noirs ».} On choisit $2$ noirs parmi $3$ et
$1$ blanc parmi $2$ : $\combi{2}{3}\times\combi{1}{2} = 3\times2 = 6$, donc
$\proba{B} = \frac{6}{10} = \frac35$.

\textbf{$C$ : « au moins un blanc ».} Le contraire est « aucun blanc »,
c'est-à-dire $A$ : $\proba{C} = 1-\frac{1}{10} = \frac{9}{10}$.
\end{exemple}

\begin{entrainement}[calculs de probabilités]
Une urne contient $4$ boules rouges et $6$ boules vertes, indiscernables au
toucher.

\exo[\niveau{1}] On tire une boule au hasard : quelle est la probabilité
qu'elle soit rouge ?

\exo[\niveau{2}] On tire simultanément $2$ boules : probabilité d'obtenir
deux rouges.

\exo[\niveau{2}] Même tirage : probabilité d'obtenir une rouge et une verte.

\exo[\niveau{2}] Même tirage : probabilité d'obtenir au moins une rouge.

\exo[\niveau{3}] On tire successivement et sans remise $3$ boules :
probabilité d'obtenir trois vertes.

\exo[\niveau{3}] On tire successivement et avec remise $3$ boules :
probabilité d'obtenir, dans l'ordre, rouge puis verte puis rouge.
\end{entrainement}

\begin{corrige}[calculs de probabilités]
\textbf{1.} $\frac{4}{10} = \frac25$.
\textbf{2.} $\frac{\combi{2}{4}}{\combi{2}{10}} = \frac{6}{45}
= \frac{2}{15}$.
\textbf{3.} $\frac{4\times6}{45} = \frac{24}{45} = \frac{8}{15}$.
\textbf{4.} Contraire : deux vertes, soit $\frac{\combi{2}{6}}{45}
= \frac{15}{45} = \frac13$ ; d'où $1-\frac13 = \frac23$.
\textbf{5.} $\frac{\arrang{3}{6}}{\arrang{3}{10}} = \frac{120}{720}
= \frac16$.
\textbf{6.} $\frac{4\times6\times4}{10^{3}} = \frac{96}{1\,000}
= \frac{12}{125}$.
\end{corrige}

\begin{qcm}
\exo Si $\proba{A} = \frac25$, alors $\proba{\evtc{A}}$ vaut :
\textbf{a)} $\frac52$ \quad \textbf{b)} $\frac35$ \quad
\textbf{c)} $-\frac25$ \quad \textbf{d)} $\frac25$ \reponseqcm{b}

\exo Une probabilité ne peut jamais valoir :
\textbf{a)} $0$ \quad \textbf{b)} $1$ \quad \textbf{c)} $\frac{4}{3}$ \quad
\textbf{d)} $0{,}99$ \reponseqcm{c}

\exo « Au moins une boule blanche » a pour contraire :
\textbf{a)} « toutes les boules sont blanches » \quad
\textbf{b)} « aucune boule blanche » \quad
\textbf{c)} « exactement une blanche » \quad
\textbf{d)} « au plus une blanche » \reponseqcm{b}

\exo Deux événements incompatibles vérifient :
\textbf{a)} $\proba{A \cap B} = \proba{A}\proba{B}$ \quad
\textbf{b)} $\proba{A \cap B} = 0$ \quad
\textbf{c)} $\proba{A \cup B} = 1$ \quad
\textbf{d)} $\proba{A} = \proba{B}$ \reponseqcm{b}

\exo Dans une urne de $10$ jetons dont $4$ rouges, on en tire $2$
simultanément. Le nombre de cas favorables à « deux rouges » est :
\textbf{a)} $\combi{2}{10}$ \quad \textbf{b)} $\combi{2}{4}$ \quad
\textbf{c)} $\arrang{2}{4}$ \quad \textbf{d)} $4^{2}$ \reponseqcm{b}
\end{qcm}

% ===========================================================================
\begin{annales}
Voici, en entier, six exercices de probabilité réellement posés. Ils
suffisent à couvrir tout ce que la série demande.

\exoann{Baccalauréat, Madagascar, série A, session 2021 — exercice 2}
Une boîte contient $5$ jetons indiscernables au toucher dont $3$ noirs et
$2$ blancs. Les résultats seront donnés sous forme de fractions
irréductibles.
\begin{enumerate}[label=\textbf{\arabic*)}, leftmargin=*, itemsep=0.15em]
  \item On tire au hasard et simultanément $3$ jetons de la boîte. Calculer
        la probabilité de chacun des événements suivants :
        $A$ : « obtenir $3$ jetons de même couleur » ;
        $B$ : « obtenir exactement deux jetons noirs » ;
        $C$ : « obtenir au moins un jeton blanc ».
  \item On tire successivement et sans remise $3$ jetons de la boîte.
        Calculer la probabilité de chacun des événements suivants :
        $D$ : « obtenir deux jetons noirs aux deux premiers tirages et un
        blanc au dernier » ;
        $E$ : « obtenir exactement un jeton blanc ».
\end{enumerate}

\exoann{Baccalauréat, Madagascar, série A, session 2022 — exercice 2}
Une trousse contient $9$ stylos indiscernables au toucher dont $5$ bleus et
$4$ rouges. Les résultats seront donnés sous forme de fraction
irréductible.
\begin{enumerate}[label=\textbf{\arabic*)}, leftmargin=*, itemsep=0.15em]
  \item On tire au hasard et simultanément $4$ stylos de la trousse.
        Calculer la probabilité de :
        $A$ : « obtenir $2$ stylos bleus et $2$ stylos rouges » ;
        $B$ : « obtenir au plus un stylo rouge ».
  \item On tire au hasard et successivement avec remise $3$ stylos.
        \begin{enumerate}[label=\textbf{\alph*)}, leftmargin=1.4em]
          \item Quel est le nombre de tirages possibles ?
          \item Calculer la probabilité de :
                $C$ : « obtenir dans l'ordre un stylo bleu et $2$ stylos
                rouges » ;
                $D$ : « obtenir au moins $2$ stylos bleus ».
        \end{enumerate}
\end{enumerate}

\exoann{Baccalauréat, Madagascar, série A, session 2020 — exercice 2}
Un sac à main contient $12$ billets de banque dont $6$ billets de
$5\,000$ Ar, $4$ billets de $10\,000$ Ar et $2$ billets de $20\,000$ Ar. On
suppose que tous les billets ont la même probabilité d'être tirés. Les
résultats seront donnés sous forme de fraction irréductible.
\begin{enumerate}[label=\textbf{\arabic*)}, leftmargin=*, itemsep=0.15em]
  \item On tire au hasard simultanément $4$ billets du sac à main.
        \begin{enumerate}[label=\textbf{\alph*)}, leftmargin=1.4em]
          \item Calculer le nombre de cas possibles.
          \item Déterminer la probabilité de :
                $A$ : « avoir exactement deux billets de $10\,000$ Ar » ;
                $B$ : « avoir quatre billets de même valeur ».
        \end{enumerate}
  \item On remet le sac dans sa condition initiale. On tire au hasard
        successivement et sans remise $3$ billets. Calculer la probabilité
        de :
        $C$ : « obtenir un billet de $20\,000$ Ar au premier tirage et
        $2$ billets de $5\,000$ Ar aux deux derniers tirages » ;
        $D$ : « obtenir un montant total de $45\,000$ Ar ».
\end{enumerate}

\exoann{Baccalauréat, Madagascar, série A, session 2004 — exercice 2}
Une trousse contient douze stylos de même marque, indiscernables au
toucher : $3$ rouges, $4$ verts, $3$ bleus et $2$ noirs. On donnera les
résultats sous forme de fraction irréductible.
\begin{enumerate}[label=\textbf{\arabic*)}, leftmargin=*, itemsep=0.15em]
  \item Un élève prend au hasard un stylo. Calculer la probabilité de :
        $A$ : « obtenir un stylo rouge » ; $B$ : « obtenir un stylo vert ou
        un stylo noir ».
  \item Un autre élève tire au hasard et simultanément deux stylos.
        Évaluer la probabilité de :
        $C$ : « les deux stylos tirés sont de même couleur » ;
        $D$ : « obtenir deux stylos de couleurs différentes ».
  \item Un troisième élève tire successivement et sans remise trois stylos.
        \begin{enumerate}[label=\textbf{\alph*)}, leftmargin=1.4em]
          \item Déterminer le nombre de cas possibles.
          \item Calculer la probabilité de :
                $E$ : « obtenir trois stylos de même couleur » ;
                $F$ : « n'obtenir aucun stylo vert ».
        \end{enumerate}
\end{enumerate}

\exoann{Baccalauréat, Madagascar, série A, session 2012 — exercice 2}
Un porte-monnaie contient $12$ billets dont $5$ billets de $500$ Ar, $4$ de
$1\,000$ Ar et $3$ de $2\,000$ Ar.
\begin{enumerate}[label=\textbf{\arabic*)}, leftmargin=*, itemsep=0.15em]
  \item On prend successivement au hasard et sans remise $3$ billets.
        \begin{enumerate}[label=\textbf{\alph*)}, leftmargin=1.4em]
          \item Vérifier qu'il y a $1\,320$ cas possibles.
          \item Calculer la probabilité d'obtenir :
                $A$ : « trois billets de même valeur » ;
                $B$ : « un montant total de $4\,500$ Ar ».
        \end{enumerate}
  \item On prend au hasard et simultanément $2$ billets. Calculer la
        probabilité d'avoir :
        $C$ : « exactement deux billets de $500$ Ar » ;
        $D$ : « au moins un billet de $1\,000$ Ar ».
\end{enumerate}

\exoann{Baccalauréat, Madagascar, série A, session 2015 — exercice 2}
On dispose de huit plaquettes indiscernables au toucher, portant chacune une
lettre du mot « SCIENCES ». On donnera les résultats sous forme de fraction
irréductible.
\begin{enumerate}[label=\textbf{\arabic*)}, leftmargin=*, itemsep=0.15em]
  \item Première épreuve : tirer simultanément quatre plaquettes.
        \begin{enumerate}[label=\textbf{\alph*)}, leftmargin=1.4em]
          \item Déterminer le nombre de cas possibles.
          \item Calculer la probabilité d'obtenir : $A$ : « aucune lettre
                E » ; $B$ : « au plus une consonne ».
        \end{enumerate}
  \item Deuxième épreuve : tirer une à une et sans remise trois plaquettes.
        Calculer la probabilité de : $C$ : « avoir au moins une voyelle » ;
        $D$ : « avoir le mot ENS ».
\end{enumerate}
\end{annales}

\begin{corrige}[annales]
\textbf{1. (2021)} Cas possibles : $\combi{3}{5} = 10$.
$\proba{A} = \frac{\combi{3}{3}}{10} = \frac{1}{10}$ ;
$\proba{B} = \frac{\combi{2}{3}\times\combi{1}{2}}{10} = \frac{6}{10}
= \frac35$ ;
$\proba{C} = 1-\proba{A} = \frac{9}{10}$.
Pour le tirage successif sans remise : $\arrang{3}{5} = 60$ cas.
$D$ : $3\times2\times2 = 12$ cas, donc $\proba{D} = \frac{12}{60}
= \frac15$.
$E$ : le blanc peut occuper l'une des trois places ($3$ choix), il y a $2$
blancs possibles, et les deux noirs se placent de $3\times2 = 6$ façons :
$3\times2\times6 = 36$, donc $\proba{E} = \frac{36}{60} = \frac35$.

\textbf{2. (2022)} Cas possibles : $\combi{4}{9} = 126$.
$\proba{A} = \frac{\combi{2}{5}\times\combi{2}{4}}{126}
= \frac{10\times6}{126} = \frac{60}{126} = \frac{10}{21}$.
$B$ : « aucun rouge » donne $\combi{4}{5} = 5$ cas, « exactement un rouge »
donne $\combi{1}{4}\times\combi{3}{5} = 4\times10 = 40$ ; total $45$, donc
$\proba{B} = \frac{45}{126} = \frac{5}{14}$.
Tirages avec remise : $9^{3} = 729$.
$\proba{C} = \frac{5\times4\times4}{729} = \frac{80}{729}$.
$D$ : exactement deux bleus, $3\times5\times5\times4 = 300$ cas ; trois
bleus, $5^{3} = 125$ cas ; total $425$, donc
$\proba{D} = \frac{425}{729}$.

\textbf{3. (2020)} Cas possibles : $\combi{4}{12} = 495$.
$\proba{A} = \frac{\combi{2}{4}\times\combi{2}{8}}{495}
= \frac{6\times28}{495} = \frac{168}{495} = \frac{56}{165}$.
$B$ : quatre billets de $5\,000$, $\combi{4}{6} = 15$ ; quatre de
$10\,000$, $\combi{4}{4} = 1$ ; impossible pour $20\,000$. Donc
$\proba{B} = \frac{16}{495}$.
Tirage successif sans remise : $\arrang{3}{12} = 1\,320$ cas.
$\proba{C} = \frac{2\times6\times5}{1\,320} = \frac{60}{1\,320}
= \frac{1}{22}$.
$D$ : le seul montant possible est $5\,000+20\,000+20\,000$ ; le billet de
$5\,000$ peut occuper l'une des trois places, d'où
$3\times6\times2\times1 = 36$ cas et
$\proba{D} = \frac{36}{1\,320} = \frac{3}{110}$.

\textbf{4. (2004)} $\proba{A} = \frac{3}{12} = \frac14$ et
$\proba{B} = \frac{4+2}{12} = \frac12$.
Tirage simultané de deux stylos : $\combi{2}{12} = 66$ cas.
$C$ : $\combi{2}{3}+\combi{2}{4}+\combi{2}{3}+\combi{2}{2}
= 3+6+3+1 = 13$, donc $\proba{C} = \frac{13}{66}$ ; et
$\proba{D} = 1-\frac{13}{66} = \frac{53}{66}$.
Tirage successif sans remise de trois stylos :
$\arrang{3}{12} = 1\,320$ cas.
$E$ : $\arrang{3}{3}+\arrang{3}{4}+\arrang{3}{3} = 6+24+6 = 36$, donc
$\proba{E} = \frac{36}{1\,320} = \frac{3}{110}$.
$F$ : on tire parmi les $8$ stylos non verts, $\arrang{3}{8} = 336$, donc
$\proba{F} = \frac{336}{1\,320} = \frac{14}{55}$.

\textbf{5. (2012)} $\arrang{3}{12} = 12\times11\times10 = 1\,320$.
\checkmark
$A$ : $\arrang{3}{5}+\arrang{3}{4}+\arrang{3}{3} = 60+24+6 = 90$, donc
$\proba{A} = \frac{90}{1\,320} = \frac{3}{44}$.
$B$ : $4\,500 = 500+2\,000+2\,000$ est la seule décomposition possible ; le
billet de $500$ occupe l'une des trois places, d'où
$3\times5\times3\times2 = 90$ cas et $\proba{B} = \frac{3}{44}$.
Tirage simultané de deux billets : $\combi{2}{12} = 66$ cas.
$\proba{C} = \frac{\combi{2}{5}}{66} = \frac{10}{66} = \frac{5}{33}$.
$D$ : le contraire est « aucun billet de $1\,000$ », soit
$\combi{2}{8} = 28$ cas ; donc
$\proba{D} = 1-\frac{28}{66} = \frac{38}{66} = \frac{19}{33}$.

\textbf{6. (2015)} Le mot SCIENCES comporte huit lettres : trois voyelles
(I, E, E) et cinq consonnes (S, S, C, C, N). Cas possibles :
$\combi{4}{8} = 70$.
$A$ : on tire quatre plaquettes parmi les six qui ne portent pas E, soit
$\combi{4}{6} = 15$ ; donc $\proba{A} = \frac{15}{70} = \frac{3}{14}$.
$B$ : « au plus une consonne » impose au moins trois voyelles ; comme il n'y
en a que trois, le seul cas est trois voyelles et une consonne :
$\combi{3}{3}\times\combi{1}{5} = 5$, donc
$\proba{B} = \frac{5}{70} = \frac{1}{14}$.
Tirage successif sans remise de trois plaquettes :
$\arrang{3}{8} = 336$ cas.
$C$ : le contraire est « que des consonnes », soit $\arrang{3}{5} = 60$
cas ; donc $\proba{C} = 1-\frac{60}{336} = \frac{276}{336} = \frac{23}{28}$.
$D$ : il faut E, puis N, puis S dans cet ordre : $2\times1\times2 = 4$ cas,
donc $\proba{D} = \frac{4}{336} = \frac{1}{84}$.
\end{corrige}

% ===========================================================================
\begin{vraifaux}
\exo Si deux événements ont la même probabilité, ils sont égaux.

\exo $\proba{A \cup B} = \proba{A}+\proba{B}$ est vrai pour tous les
événements.

\exo Une probabilité nulle signifie que l'événement est impossible dans un
univers fini équiprobable.

\exo Dans un tirage simultané, l'ordre des objets tirés compte.

\exo Le contraire de « exactement deux rouges » est « au moins trois
rouges ».
\end{vraifaux}

\begin{corrige}[vrai ou faux]
\textbf{1.} Faux — obtenir pile et obtenir face ont la même probabilité sans
être le même événement.
\textbf{2.} Faux — seulement si $A$ et $B$ sont incompatibles.
\textbf{3.} Vrai — aucun cas favorable, donc l'événement ne peut pas se
produire.
\textbf{4.} Faux — c'est précisément ce qui le distingue du tirage
successif.
\textbf{5.} Faux — le contraire est « zéro, une, ou au moins trois
rouges », autrement dit « pas exactement deux ».
\end{corrige}

% ===========================================================================
\begin{situation}[la tombola du collège]
Pour financer un voyage, un collège vend $200$ billets de tombola numérotés
de $1$ à $200$. Trois billets seront tirés au sort, simultanément, pour
désigner les trois gagnants. Une classe a acheté $12$ billets.

\begin{enumerate}[label=\textbf{\arabic*.}, leftmargin=*, itemsep=0.3em]
  \item Combien y a-t-il de tirages possibles ? (On laissera le résultat
        sous forme de combinaison.)
  \item Exprimer, sous forme de quotient de combinaisons, la probabilité que
        les trois gagnants appartiennent à cette classe.
  \item Exprimer de même la probabilité qu'aucun des trois gagnants
        n'appartienne à cette classe.
  \item En déduire la probabilité qu'au moins un gagnant appartienne à la
        classe, puis en donner une valeur approchée à $10^{-2}$ près.
\end{enumerate}
\end{situation}

\begin{corrige}[la tombola du collège]
\textbf{1.} $\combi{3}{200}$ tirages possibles, soit
$\frac{200\times199\times198}{6} = 1\,313\,400$.

\textbf{2.} $\dfrac{\combi{3}{12}}{\combi{3}{200}}
= \dfrac{220}{1\,313\,400} \approx 0{,}000\,17$.

\textbf{3.} $\dfrac{\combi{3}{188}}{\combi{3}{200}}
= \dfrac{1\,085\,846}{1\,313\,400}$.

\textbf{4.} La probabilité cherchée vaut
$1-\dfrac{\combi{3}{188}}{\combi{3}{200}}
\approx 1-0{,}827 = 0{,}173$, soit environ $0{,}17$ : à peu près une chance
sur six.
\end{corrige}

% ===========================================================================
\begin{jemeteste}
Vingt-cinq minutes. Une urne contient $7$ jetons : $4$ verts et $3$ jaunes.

\begin{enumerate}[label=\textbf{\arabic*.}, leftmargin=*, itemsep=0.3em]
  \item On tire simultanément $2$ jetons : combien de cas possibles ?
  \item Probabilité d'obtenir deux jetons verts.
  \item Probabilité d'obtenir deux jetons de couleurs différentes.
  \item Probabilité d'obtenir au moins un jaune.
  \item On tire successivement et avec remise $3$ jetons : probabilité
        d'obtenir, dans l'ordre, vert, vert, jaune.
\end{enumerate}
\end{jemeteste}

\begin{corrige}[je me teste]
\textbf{1.} $\combi{2}{7} = 21$.
\textbf{2.} $\frac{\combi{2}{4}}{21} = \frac{6}{21} = \frac27$.
\textbf{3.} $\frac{4\times3}{21} = \frac{12}{21} = \frac47$.
\textbf{4.} Contraire : deux verts, donc $1-\frac27 = \frac57$.
\textbf{5.} $\frac{4\times4\times3}{7^{3}} = \frac{48}{343}$.
\end{corrige}

% ===========================================================================
\begin{commeaubac}
\bmtsource{Baccalauréat, Madagascar, série A, session 2016 — exercice 2
\textbullet\ 5 points}
Une boîte contient $10$ crayons indiscernables au toucher dont
\begin{itemize}[leftmargin=1.6em, itemsep=0.1em]
  \item $3$ rouges numérotés $1$ ; $2$ ; $3$,
  \item $3$ blancs numérotés $1$ ; $1$ ; $2$,
  \item $4$ verts numérotés $1$ ; $2$ ; $4$ ; $4$.
\end{itemize}
Chaque crayon a la même probabilité d'être tiré. On donnera les résultats
sous forme de fraction irréductible.

\begin{enumerate}[label=\textbf{\arabic*.}, leftmargin=*, itemsep=0.25em]
  \item Un enfant prend, au hasard et d'un seul coup, $3$ crayons de la
        boîte. Calculer la probabilité de chacun des événements suivants :
    \begin{enumerate}[label=\textbf{\alph*)}, leftmargin=1.4em, itemsep=0.12em]
      \item $A$ : « obtenir $3$ crayons de couleurs différentes » ;
            \bareme{1 pt}
      \item $B$ : « obtenir exactement $2$ crayons blancs ». \bareme{1 pt}
    \end{enumerate}
  \item Un autre enfant prend au hasard, un à un avec remise, $3$ crayons de
        la boîte. Calculer la probabilité de chacun des événements
        suivants :
    \begin{enumerate}[label=\textbf{\alph*)}, leftmargin=1.4em, itemsep=0.12em]
      \item $C$ : « obtenir, dans l'ordre, un crayon rouge et deux crayons
            verts » ; \bareme{1 pt}
      \item $D$ : « obtenir $3$ crayons dont la somme des numéros est égale
            à $5$ » ; \bareme{1 pt}
      \item $E$ : « obtenir $3$ crayons dont le produit des numéros est égal
            à $8$ ». \bareme{1 pt}
    \end{enumerate}
\end{enumerate}
\end{commeaubac}

\begin{corrige}[comme au baccalauréat — Madagascar A 2016]
\textbf{1.} Cas possibles : $\combi{3}{10} = 120$.

\textbf{a)} Un crayon de chaque couleur : $3\times3\times4 = 36$ cas, donc
$\proba{A} = \frac{36}{120} = \frac{3}{10}$.

\textbf{b)} Deux blancs parmi trois et un crayon parmi les sept autres :
$\combi{2}{3}\times\combi{1}{7} = 3\times7 = 21$, donc
$\proba{B} = \frac{21}{120} = \frac{7}{40}$.

\textbf{2.} Cas possibles : $10^{3} = 1\,000$.

\textbf{a)} Dans l'ordre : un rouge ($3$ choix), puis deux verts ($4$ puis
$4$) : $3\times4\times4 = 48$ cas, donc
$\proba{C} = \frac{48}{1\,000} = \frac{6}{125}$.

\textbf{b)} Comptons d'abord les crayons par numéro : le $1$ figure sur
quatre crayons (un rouge, deux blancs, un vert), le $2$ sur trois, le $3$
sur un, le $4$ sur deux. Les sommes égales à $5$ s'obtiennent avec les
numéros $\{1\,;\,1\,;\,3\}$ ou $\{1\,;\,2\,;\,2\}$.
Pour $(1,1,3)$ : trois places possibles pour le $3$, et
$4\times4\times1 = 16$ choix de crayons, soit $3\times16 = 48$ cas.
Pour $(1,2,2)$ : trois places possibles pour le $1$, et
$4\times3\times3 = 36$ choix, soit $3\times36 = 108$ cas.
Total : $156$ cas, donc $\proba{D} = \frac{156}{1\,000} = \frac{39}{250}$.

\textbf{c)} Le produit vaut $8$ pour les numéros $\{1\,;\,2\,;\,4\}$ ou
$\{2\,;\,2\,;\,2\}$.
Pour $(1,2,4)$ : $3! = 6$ ordres, et $4\times3\times2 = 24$ choix, soit
$144$ cas. Pour $(2,2,2)$ : $3^{3} = 27$ cas.
Total : $171$ cas, donc $\proba{E} = \frac{171}{1\,000}$.
\end{corrige}

\begin{notepourlaclasse}
Huit heures, le plus gros volume horaire du livre après le logarithme, et
c'est justifié : l'exercice de probabilité est celui où les élèves de la
série marquent le plus de points quand il est travaillé, et le moins quand
il ne l'est pas.

Faites écrire, pour chaque question, les trois lignes rituelles : le nombre
de cas possibles, le nombre de cas favorables, la fraction simplifiée. Les
correcteurs les cherchent, et elles se notent séparément.

La question 2.b) du sujet 2016 — la somme des numéros — est la plus
difficile de tout le fonds : elle demande de trier les crayons par numéro et
non par couleur. Traitez-la au tableau, avec le tableau des effectifs par
numéro construit devant les élèves ; sans lui, elle est hors de portée.

Rappelez enfin, à chaque séance, que les probabilités conditionnelles et les
variables aléatoires ne sont pas au programme de la série : des élèves s'y
égarent chaque année en consultant un autre manuel, sur des notions qu'on ne
leur demandera pas.
\end{notepourlaclasse}
